%! The Complete Solution to Ax = b — Unit 3, Lesson 3.3 — Group Activity
\documentclass[10pt]{article}
\usepackage{linalg-article}
\usepackage{linalg-boxes}
\newcommand{\TallMath}[1]{$\displaystyle #1\rule[-1.4em]{0pt}{3.2em}$}
\newcommand{\vv}{\mathbf{v}}
\newcommand{\ww}{\mathbf{w}}
\newcommand{\xx}{\mathbf{x}}
\newcommand{\yy}{\mathbf{y}}
\newcommand{\bb}{\mathbf{b}}
\newcommand{\svec}{\mathbf{s}}
\newcommand{\zero}{\mathbf{0}}

\begin{document}

\pageheader{Unit 3, Lesson 3.3}{Group Activity}
\namepartnerperiod

\vspace{0.10in}

\begin{headlinebox}{blush}
\small\textbf{The Complete Solution.} Every solution of $A\xx=\bb$ is $\xx=\xx_p+\xx_n$: \emph{one}
particular solution plus the \emph{whole} nullspace. Find $\xx_p$ by reducing $[A\mid\bb]$ and setting
the free variables to $0$; attach the special solutions for the nullspace part. The solution set is
$N(A)$ \emph{shifted} to pass through $\xx_p$ --- and $A\xx=\bb$ is solvable only if $\bb$ lies in $C(A)$.
Work with your partner and \emph{justify} every answer.
\end{headlinebox}

\vspace{0.10in}

\begin{tcolorbox}[colback=white, colframe=black!40, title=\textbf{Tier R --- Remediate},
    fonttitle=\bfseries, arc=1.5mm, left=3mm, right=3mm, top=2mm, bottom=2mm, breakable]
Build a complete solution from its two pieces. Use $A=\begin{bmatrix} 1 & 2 \\ 2 & 4 \end{bmatrix}$ and
$\bb=\begin{bmatrix} 3 \\ 6 \end{bmatrix}$ throughout.
\begin{enumerate}[leftmargin=1.7em, itemsep=8pt]
  \item \textbf{Check the particular solution.} Verify $\xx_p=\begin{bsmallmatrix} 3\\0 \end{bsmallmatrix}$ solves $A\xx=\bb$:
        $A\xx_p=\begin{bmatrix}\rule{0pt}{1.0em}\blank{1.0cm}\\[2pt]\blank{1.0cm}\end{bmatrix}$. Does it equal $\bb$? \; \blank{1.4cm}
  \item \textbf{Add a do-nothing move.} The nullspace is all multiples of $\svec=\begin{bsmallmatrix} -2\\1 \end{bsmallmatrix}$.
        Check that $\xx_p+\svec=\begin{bsmallmatrix} 1\\1 \end{bsmallmatrix}$ \emph{also} solves $A\xx=\bb$:
        $A(\xx_p+\svec)=\begin{bmatrix}\rule{0pt}{1.0em}\blank{1.0cm}\\[2pt]\blank{1.0cm}\end{bmatrix}$. \; \blank{1.4cm}
  \item \textbf{Write the complete solution} by combining the two pieces:
        $\xx=\blank{1.6cm}+\;t\,\blank{1.6cm}$.
\end{enumerate}
\end{tcolorbox}

\begin{tcolorbox}[colback=white, colframe=black!40, title=\textbf{Tier A --- Approaching Proficiency},
    fonttitle=\bfseries, arc=1.5mm, left=3mm, right=3mm, top=2mm, bottom=2mm, breakable]
For each system: reduce $[A\mid\bb]$, set the free variable(s) to $0$ to read $\xx_p$, attach the
nullspace, write the complete solution, and say what shape the solution set is.
\begin{enumerate}[leftmargin=1.7em, itemsep=9pt]
  \item $A=\begin{bmatrix} 1 & 2 & 1 \\ 2 & 4 & 3 \end{bmatrix}$, $\bb=\begin{bmatrix} 4 \\ 9 \end{bmatrix}$ \quad (\emph{the free column is the middle one}). \writelines{3}
  \item $A=\begin{bmatrix} 1 & 2 & 2 \\ 2 & 4 & 4 \end{bmatrix}$, $\bb=\begin{bmatrix} 3 \\ 6 \end{bmatrix}$ \quad (\emph{expect two free variables}). \writelines{3}
  \item \textbf{Through the origin?} Is the solution set in problem~2 a subspace --- does it pass through
        $\zero$? Explain in one sentence. \writelines{2}
\end{enumerate}
\end{tcolorbox}

\begin{tcolorbox}[colback=white, colframe=black!40, title=\textbf{Tier E --- Extension},
    fonttitle=\bfseries, arc=1.5mm, left=3mm, right=3mm, top=2mm, bottom=2mm, breakable]
\textbf{Solvability, uniqueness, and a proof.}
\begin{enumerate}[leftmargin=1.7em, itemsep=9pt]
  \item \textbf{Which $\bb$ are reachable?} For $A=\begin{bmatrix} 1 & 2 \\ 3 & 6 \end{bmatrix}$, the columns
        both lie on the line $y=3x$. For which right-hand sides $\bb=\begin{bsmallmatrix} b_1\\b_2 \end{bsmallmatrix}$
        is $A\xx=\bb$ solvable? Give one solvable $\bb$ and one that is not. \writelines{2}
  \item \textbf{When is the solution unique?} Explain why a \emph{solvable} system $A\xx=\bb$ has exactly
        one solution precisely when $N(A)=\{\zero\}$ (no free variables). \writelines{2}
  \item \textbf{Prove the form.} Suppose $\xx$ and $\xx_p$ both solve $A\xx=\bb$. Show that $\xx-\xx_p$ is
        in $N(A)$, and conclude that \emph{every} solution has the form $\xx_p+\xx_n$. \writelines{3}
\end{enumerate}
\end{tcolorbox}

\end{document}
